\section{Conclusion and Limitations}
\label{sec:conclusion}

VLA quantization is a closed-loop, distribution-dependent problem: quantized actions change future observations and hence the data used to judge the policy. \method restores temporal context with a physical-sequence contract, whole-policy context with an abstaining exact-budget audit, and deployment context with input-conditioned activation ranges, using one static W4/FP16 mask without task routing or output correction.

On GR00T, this yields 54.0\% RoboCasa365 success at 2.22$\times$ static compression, versus 30.4\% for size-matched \quantvla W4A8 and 55.1\% for FP16; the latter comparison does not establish equivalence. In a paired 500-episode intervention, the registered frozen A8 table obtains 0.0\% and input-conditioned A8 53.8\%. The GR00T audit abstains, and alternative masks isolate no separate mean-success gain.

The protocol-matched, four-flow-step $\pi_{0.5}$ point is 27.7\% at 2.702$\times$ compression, but its gaps and mask change remain descriptive. Missing tests include teacher-forced feedback isolation, \dpac predictive validation, a closed-loop FP16-activation arm, FP16 non-inferiority, and real latency, memory, and physical-safety measurements. They bound the present claim and define the next experiments.
